\clearpage
\fancyhead[L]{\small SONNENPARALLAXE.}
\SourceLeaf{528}
\begin{center}\small Göttingische gelehrte Anzeigen. Stück 17. S. 161..166. 29. Januar 1810.\\[3pt]\rule{.72\linewidth}{.3pt}\end{center}

Herr \textsc{Lars Regner}, Professor der Astronomie in Uppsala und Mitglied der dortigen Königlichen Gesellschaft der Wissenschaften, hat unter dem Titel \emph{Supplementum ad historiam de Parallaxeos Solaris inventione, illustriss.\ scientiarum Societatis regiae Gottingensis dijudicationi subjectum} einen Aufsatz eingesandt, dessen wesentlicher Inhalt im Folgenden wiedergegeben wird. Herr \textsc{Regner} erzählt zunächst, er sei schon 1806 davon überzeugt gewesen, dass die von den Astronomen angenommene Sonnenparallaxe unrichtig sei. Anfang des folgenden Jahres habe er eine kleine Schrift darüber an die Petersburger Akademie der Wissenschaften gesandt. Herr \textsc{Fuss} habe als Sekretär dieser Akademie geantwortet, die Abhandlung sei Herrn \textsc{Schubert} zur Beurteilung übergeben worden, der sie für irrig und falsch befunden habe. Er selbst zweifle auch gar nicht daran, dass es ausreichen werde, Herrn \textsc{Regner} die Ursache der Fehler lediglich anzudeuten, um ihn davon zu überzeugen. Überdies biete er an, ihm den ganzen Bericht selbst mitzuteilen. Herr \textsc{Regner} wartete diesen aber nicht ab. In der Überzeugung, dass der angedeutete Tadel ungerecht sei, ließ er seine Schrift mit einigen Erläuterungen in Uppsala drucken und schickte sie dann erneut nach Petersburg. Daraufhin sandte ihm Herr \textsc{Fuss} ohne Weiteres den Bericht des Herrn \textsc{Schubert} selbst.

Herr \textsc{Regner} lässt nun einen Auszug aus seiner gedruckten Schrift folgen (die uns sonst nicht bekannt geworden ist). Das Wesentliche daran ist die Bemerkung, dass die Anziehung des Mondes durch die Sonne beträchtlich größer als die Anziehung des Mondes durch die Erde ist, wenn man die Sonnenparallaxe zwischen $8''5$ und $8''9$ annimmt. Aus dieser (allerdings richtigen) Bemerkung glaubt Herr \textsc{Regner} schließen zu müssen, dass jene Parallaxe unmöglich richtig sein könne. Denn der Mond würde, da ihn die Sonne stärker anzieht als die Erde, (nach der Vorstellung des Herrn \textsc{Regner}) die Erde verlassen und seine eigene Bahn um die Sonne beschreiben müssen. Ja, selbst wenn man nur annähme, dass die Sonne den Mond ebenso stark anzieht wie die Erde, würden bei Neumond beide Anziehungen einander vollständig aufheben. Der Mond würde sich entlang der Tangente weiterbewegen und dadurch sofort in einen Bereich geraten, in dem die Anziehung der Sonne überwiegt. Diese letzte Behauptung will Herr \textsc{Regner} durch eine Figur verdeutlichen, in der er die Erde, den Mond und die Sonne auf einer Geraden mit $T$, $b$, $S$ bezeichnet. Dann zeigt er, dass sich die Punkte der in $b$ errichteten senkrechten Geraden von $T$ schneller entfernen würden als von $S$. (Wie konnte Herr Professor \textsc{Regner} vergessen, dass die Erde, während sich der Mond entlang der Tangente weiterbewegte, ja nicht in $T$ bleibt, sondern sich aufgrund ihrer Tangentialgeschwindigkeit von selbst weiterbewegt, sodass der Abstand des Mondes von ihr keineswegs abnehmen muss? Ebenso liegt in einer \emph{größeren} Anziehung zur Sonne hin nichts Widersprüchliches. Es folgt \emph{nur}, dass die Mondbahn im absoluten Raum bei Neumond zur Sonne hin konkav werden muss. Dass sie es wirklich ist, davon kann man sich leicht überzeugen, wenn man sie auch nur unter Vernachlässigung der Ungleichheiten in der Bewegung der Erde und des Mondes als Epizykloide betrachtet.)

Herr \textsc{Regner} fügt nun einen Auszug aus dem Bericht des Herrn \textsc{Schubert} bei. Herr \textsc{Regner} müsse, so urteilt \textsc{Schubert}, die Abhandlung entweder in großer Eile oder im Scherz geschrieben haben. Denn wäre sie ernst gemeint, so würde daraus folgen, dass der Verfasser weder die Methode kenne, mit der die Sonnenparallaxe bestimmt worden ist, noch die \emph{ersten Grundsätze} des Dreikörperproblems. Die erste Behauptung begründet \textsc{Schubert} damit, dass man in den Beobachtungen des Venusdurchgangs einen Fehler von 27 Minuten annehmen müsste, wenn die Sonnenparallaxe 21 Sekunden betrüge, die Herr \SourceLeaf{529}\textsc{Regner} nach seiner vorhin erwähnten Vorstellung für den kleinsten zulässigen Wert hält. Zur Bekräftigung der zweiten Behauptung zeigt \textsc{Schubert}, dass Herr \textsc{Regner} die absolute Bewegung des Mondes mit seiner Bewegung relativ zur Erde verwechselt habe. Bei der Betrachtung der letzteren ist die Anziehung des Mondes durch die Sonne für sich genommen unerheblich. Maßgeblich ist allein, wie stark sie sich von der Anziehung der Erde durch die Sonne unterscheidet. Dieser Unterschied sei aber im Vergleich zur Anziehung des Mondes durch die Erde nur gering. (Leser, denen die Grundbegriffe der physischen Astronomie nicht ganz fremd sind, werden uns eine weitere Erörterung dieser Ausführungen von \textsc{Schubert} gern erlassen.)

Hierauf lässt Herr \textsc{Regner} schließlich Animadversiones in relationem D\textsuperscript{nı} \textsc{Schubert} ad Scient. Imperial. Petropolitanam de Scripto parallaxin Solis tractante folgen. Wir übergehen die langen Klagen über die Ungerechtigkeit, durch die er sich von Herrn \textsc{Schubert} gekränkt und beschimpft glaubt, und geben nur noch kurz an, wie Herr \textsc{Regner} die Beschuldigungen von sich abzuwälzen sucht. Zunächst bemüht er sich, die Genauigkeit der Beobachtungen bei den Venusdurchgängen von 1761 und 1769 in Zweifel zu ziehen, die \textsc{Lalande} auf ein paar Zeitsekunden oder ein paar Zehntel einer Bogensekunde angesetzt hatte. Eine solche Genauigkeit sei mit den damals verwendeten Fernrohren schlechthin unerreichbar gewesen. \emph{Deswegen} müsse man eine Anhäufung von Beobachtungsfehlern bis auf 27 Zeitminuten durchaus für möglich halten und die damaligen Beobachtungen für die Bestimmung der Sonnenparallaxe für völlig untauglich erklären! (Wir können uns unmöglich entschließen, über eine solche Art der Rechtfertigung noch etwas hinzuzufügen.)

Herr \textsc{Regner} glaubt nun, ihm bleibe nichts weiter übrig, als noch die zweite Beschuldigung des Herrn \textsc{Schubert} zu entkräften. Er beginnt mit einigen Gemeinplätzen über die Nachteile blinder Unterwerfung unter Autorität und überlieferte Irrtümer in den Wissenschaften, die seiner Meinung nach gerade hier am Platz seien. Allein auf \textsc{Newton}s Wort hin glaube man seither, dass die Wirkung der Anziehung durch die Sonne auf Erde und Mond in gleichen Abständen gleich groß sei. Dieser Satz sei jedoch völlig falsch, und Herr \textsc{Regner} setzt an seine Stelle folgenden Satz. Wirken gleiche Kräfte auf zwei Körper mit ungleichen Massen, so erhalten diese ungleiche Geschwindigkeiten, die zu den Massen umgekehrt proportional sind. (Warum muss man Herrn Professor \textsc{Regner} hier daran erinnern, dass dieser Satz nur gelten kann, wenn von gleicher \emph{bewegender} Kraft die Rede ist, wie man sie in der Mechanik von der beschleunigenden Kraft unterscheidet, und dass die Sonne auf die Himmelskörper nur insofern bewegende Kraft ausübt, als sie auf ihre einzelnen Teile wirkt, also umso mehr bewegende Kraft auf das Ganze ausübt, je mehr Teile angezogen werden?) Um diesen Satz zu verteidigen, bestreitet Herr \textsc{Regner} geradezu die Beweiskraft der Versuche, die mit fallenden Körpern unter der Glocke der Luftpumpe angestellt worden sind. Teils seien die Fallzeiten zu kurz, teils könne man die Luft nie \emph{ganz} entfernen. Dass unter der Glocke die Feder ebenso schnell fällt wie der Dukaten, liegt nach Ansicht des Herrn \textsc{Regner} allein daran, dass die restliche Luft der Feder mehr Widerstand entgegensetzt als dem Goldstück. Im völlig leeren Raum würde die Feder um so viel schneller fallen, wie sie weniger Masse hat. Ebenso glaubt er, gerade sein Satz werde unwiderleglich bewiesen durch das Phänomen, dass zwei Kugeln mit gleichem Volumen und ungleichen Massen, an gleich langen Fäden aufgehängt, gleich schnelle Pendelschwingungen ausführen (was übrigens, sofern der Luftwiderstand noch merklich ist, gar nicht stimmt). Auf die kleinere Masse müsse nämlich eine um so viel größere Kraft wirken, damit sie denselben Widerstand ebenso leicht überwindet und ebenso schnell fällt. Zuletzt findet Herr \textsc{Regner} noch einen Beweis seines Satzes im Gleichgewicht zweier Körper ungleicher Masse an ungleichen Hebelarmen, deren Längen zu den Massen umgekehrt proportional sind, weil die virtuellen Geschwindigkeiten hier im umgekehrten Verhältnis der Massen stehen. Da Herr \textsc{Regner} aufgrund all dessen für ausgemacht hält, dass \SourceLeaf{530}die Wirkung der Anziehung durch die Sonne auf den Mond um so viel größer ist als auf die Erde, wie die Masse des Mondes kleiner ist als die der Erde, fällt es ihm leicht, die Unzulässigkeit der Schlüsse des Herrn \textsc{Schubert}s zu zeigen.

Weder der Raum noch der Zweck dieser Blätter gestatten uns, über diese Darlegung hinauszugehen und unseren Lesern die Schlüsse vorwegzunehmen, die sie daraus leicht selbst ziehen werden.
